\mxunit{premiere-progression}{Conduite annuelle et créneaux}\label{ANNUEL}
70 occurrences présentes ; 2 créneaux neutralisés pour les fenêtres de blancs ; 4 créneaux après le 18 juin conditionnels. 64 séances ordinaires avant le 18 juin, soit 128 h, dont deux entraînements complets. DS du samedi et blancs : durée et placement à confirmer, non ajoutés artificiellement au volume de cours.
\paragraph{Trame de 120 min.} Rituel 8 ; retour du travail 10 ; recherche ou reprise 15 ; cours et exemple 30 ; exercices 40 ; bilan et travail personnel 7 ; marge de régulation 10. Pour une première séance, le diagnostic prend une partie du temps d’exercices. Pour un sujet complet, les 120 min sont réservées au sujet, sans rituel ajouté.
\mxneed{8}\section*{S01 -- 2026-09-03 : Accueil et diagnostic transversal}
\textbf{Statut :} cours. 09:45–11:45.\par\textbf{Rituel :} R001, R008.\par\textbf{En classe :} Diagnostic des acquis de seconde et présentation des méthodes de travail.\par\textbf{Travail personnel :} Finir uniquement les questions non traitées du diagnostic, 10 min.\par\textbf{Note :} Acquisition prévue, à ajuster à partir du travail réellement effectué et des diagnostics.
\mxneed{8}\section*{S02 -- 2026-09-08 : Définir et représenter une suite}
\textbf{Statut :} cours. 09:45–11:45.\par\textbf{Rituel :} R002, R009.\par\textbf{En classe :} C01-U1 ; exemples résolus et applications ; C01-E1. Diagnostic D : 10 min ; activité RECH : 15 min.\par\textbf{Travail personnel :} CQFD, exercice 9 p. 64, questions a et b. Calculer les quatre premiers termes. Écrire les substitutions avant de simplifier. 10–20 min maximum ; remplace le travail personnel initial.\par\textbf{Note :} Acquisition prévue, à ajuster à partir du travail réellement effectué et des diagnostics. Correction CQFD dans le livret professeur ; reprendre au prochain cours.
\mxneed{8}\section*{S03 -- 2026-09-10 : Suites arithmétiques}
\textbf{Statut :} cours. 09:45–11:45.\par\textbf{Rituel :} R003, R031.\par\textbf{En classe :} C01-U2 ; exemples résolus et applications ; C01-E2 : modèle A seulement. \par\textbf{Travail personnel :} CQFD, exercice 11 p. 64, questions toutes. Calculer les termes successivement et conserver la trace des quatre étapes. 10–20 min maximum ; remplace le travail personnel initial.\par\textbf{Note :} Acquisition prévue, à ajuster à partir du travail réellement effectué et des diagnostics. Correction CQFD dans le livret professeur ; reprendre au prochain cours.
\mxneed{8}\section*{S04 -- 2026-09-15 : Suites géométriques}
\textbf{Statut :} cours. 09:45–11:45.\par\textbf{Rituel :} R004, R032.\par\textbf{En classe :} C01-U3 ; exemples résolus et applications ; C01-E3 ; E2 modèle B. \par\textbf{Travail personnel :} C01-E3, 10–20 min maximum.\par\textbf{Note :} Acquisition prévue, à ajuster à partir du travail réellement effectué et des diagnostics.
\mxneed{8}\section*{S05 -- 2026-09-17 : Sommes, indices et démonstrations}
\textbf{Statut :} cours. 09:45–11:45.\par\textbf{Rituel :} R005, R032.\par\textbf{En classe :} C01-U4, C01-U5 ; exemples résolus et applications ; C01-E4 ; E6 selon temps disponible.  QCM 1 C01 (10 min, correction comprise). QCM 2 C01 (10 min) et bilan.\par\textbf{Travail personnel :} C01-E8 à reprendre dans une semaine, 10–20 min maximum.\par\textbf{Note :} Acquisition prévue, à ajuster à partir du travail réellement effectué et des diagnostics. Sommes : formules et preuves essentielles ; exploration des seuils à reprendre le 15 décembre. Pas de récurrence formelle.
\mxneed{8}\section*{S06 -- 2026-09-22 : Trois formes, trois lectures}
\textbf{Statut :} cours. 09:45–11:45.\par\textbf{Rituel :} R006, R031.\par\textbf{En classe :} C02-U1 ; exemples résolus et applications ; C02-E1. Diagnostic D : 10 min ; activité RECH : 15 min.\par\textbf{Travail personnel :} C02-E1, 10–20 min maximum.\par\textbf{Note :} Acquisition prévue, à ajuster à partir du travail réellement effectué et des diagnostics.
\mxneed{8}\section*{S07 -- 2026-09-24 : Discriminant et résolution}
\textbf{Statut :} cours. 09:45–11:45.\par\textbf{Rituel :} R007, R031.\par\textbf{En classe :} C02-U2 ; exemples résolus et applications ; C02-E3. \par\textbf{Travail personnel :} CQFD, exercice 43 p. 99, questions a et c. Trouver d’abord une racine évidente, puis l’autre à l’aide du produit des racines. 10–20 min maximum ; remplace le travail personnel initial.\par\textbf{Note :} Acquisition prévue, à ajuster à partir du travail réellement effectué et des diagnostics. Correction CQFD dans le livret professeur ; reprendre au prochain cours.
\mxneed{8}\section*{S08 -- 2026-09-29 : Factoriser sans automatiser le discriminant}
\textbf{Statut :} cours. 09:45–11:45.\par\textbf{Rituel :} R008, R031.\par\textbf{En classe :} C02-U3 ; exemples résolus et applications ; C02-E4. \par\textbf{Travail personnel :} CQFD, exercice 49 p. 99, questions a et b. Traduire la somme et le produit par une équation ; ne pas supposer d’avance que les nombres existent. 10–20 min maximum ; remplace le travail personnel initial.\par\textbf{Note :} Acquisition prévue, à ajuster à partir du travail réellement effectué et des diagnostics. Correction CQFD dans le livret professeur ; reprendre au prochain cours.
\mxneed{8}\section*{S09 -- 2026-10-01 : Signe, inéquations et optimisation}
\textbf{Statut :} cours. 09:45–11:45.\par\textbf{Rituel :} R009, R031.\par\textbf{En classe :} C02-U4 ; exemples résolus et applications ; C02-E5 ou E6.  QCM 1 C02 (10 min, correction comprise). QCM 2 C02 (10 min) et bilan.\par\textbf{Travail personnel :} C02-E8 à reprendre dans une semaine, 10–20 min maximum.\par\textbf{Note :} Acquisition prévue, à ajuster à partir du travail réellement effectué et des diagnostics.
\mxneed{8}\section*{S10 -- 2026-10-06 : Projection et angle}
\textbf{Statut :} cours. 09:45–11:45.\par\textbf{Rituel :} R010, R031.\par\textbf{En classe :} C03-U1 ; exemples résolus et applications ; C03-E4 et E6. Diagnostic D : 10 min ; activité RECH : 15 min.\par\textbf{Travail personnel :} C03-E4, 10–20 min maximum.\par\textbf{Note :} Acquisition prévue, à ajuster à partir du travail réellement effectué et des diagnostics.
\mxneed{8}\section*{S11 -- 2026-10-08 : Coordonnées, norme et orthogonalité}
\textbf{Statut :} cours. 09:45–11:45.\par\textbf{Rituel :} R011, R031.\par\textbf{En classe :} C03-U2 ; exemples résolus et applications ; C03-E1 et E2. \par\textbf{Travail personnel :} CQFD, exercice 9 p. 250, questions a et b. Calculer le produit scalaire dans la base orthonormée et détailler les deux produits de coordonnées. 10–20 min maximum ; remplace le travail personnel initial.\par\textbf{Note :} Acquisition prévue, à ajuster à partir du travail réellement effectué et des diagnostics. Correction CQFD dans le livret professeur ; reprendre au prochain cours.
\mxneed{8}\section*{S12 -- 2026-10-13 : Calcul vectoriel et identités}
\textbf{Statut :} cours. 09:45–11:45.\par\textbf{Rituel :} R012, R031.\par\textbf{En classe :} C03-U3 ; exemples résolus et applications ; C03-E3.  QCM 1 C03 (10 min, correction comprise).\par\textbf{Travail personnel :} CQFD, exercice 11 p. 250, questions 1, 2 et 3. Utiliser successivement les coordonnées, la relation de Chasles et le fait que I est le milieu. 10–20 min maximum ; remplace le travail personnel initial.\par\textbf{Note :} Acquisition prévue, à ajuster à partir du travail réellement effectué et des diagnostics. Correction CQFD dans le livret professeur ; reprendre au prochain cours.
\mxneed{8}\section*{S13 -- 2026-10-15 : Al-Kashi et choix d’une méthode}
\textbf{Statut :} cours. 09:45–11:45.\par\textbf{Rituel :} R013, R031.\par\textbf{En classe :} C03-U4, C03-U5 ; exemples résolus et applications ; C03-E8 et E5 selon temps disponible.  QCM 2 C03 (10 min) et bilan.\par\textbf{Travail personnel :} C03-E8 à reprendre dans une semaine, 10–20 min maximum.\par\textbf{Note :} Acquisition prévue, à ajuster à partir du travail réellement effectué et des diagnostics.
\mxneed{8}\section*{S14 -- 2026-11-03 : Taux de variation et sécantes}
\textbf{Statut :} cours. 09:45–11:45.\par\textbf{Rituel :} R014, R038.\par\textbf{En classe :} C04-U1 ; exemples résolus et applications ; C04-E1 : calcul des deux taux seulement, comparaison au nombre dérivé différée. Diagnostic D : 10 min ; activité RECH : 15 min. Mise à disposition de EV1-T1 : choisir une question adaptée aux notions déjà étudiées.\par\textbf{Travail personnel :} C04-E1 : les taux seulement, 10–20 min maximum.\par\textbf{Note :} Acquisition prévue, à ajuster à partir du travail réellement effectué et des diagnostics. Entraînement de réserve : remplace une tâche personnelle prévue, sans cumul de devoirs ; ne pas exiger deux sujets complets à domicile.
\mxneed{8}\section*{S15 -- 2026-11-05 : Nombre dérivé}
\textbf{Statut :} cours. 09:45–11:45.\par\textbf{Rituel :} R015, R039.\par\textbf{En classe :} C04-U2 ; exemples résolus et applications ; C04-E2 et E3.  Mise à disposition de EV1-T2 : choisir une question adaptée aux notions déjà étudiées.\par\textbf{Travail personnel :} CQFD, exercice 17 p. 127, questions toutes. Lire deux points de chaque tangente ; écrire le quotient des accroissements plutôt qu’estimer son inclinaison. 10–20 min maximum ; remplace le travail personnel initial.\par\textbf{Note :} Acquisition prévue, à ajuster à partir du travail réellement effectué et des diagnostics. Entraînement de réserve : remplace une tâche personnelle prévue, sans cumul de devoirs ; ne pas exiger deux sujets complets à domicile. Correction CQFD dans le livret professeur ; reprendre au prochain cours.
\mxneed{8}\section*{S16 -- 2026-11-10 : Tangente et approximation affine}
\textbf{Statut :} cours. 09:45–11:45.\par\textbf{Rituel :} R016, R040.\par\textbf{En classe :} C04-U3 ; exemples résolus et applications ; C04-E4 et E5.  QCM 1 C04 (10 min, correction comprise).\par\textbf{Travail personnel :} CQFD, exercice 19 p. 127, questions 1 et 2. Distinguer l’ordonnée du point de contact et la pente de la droite tangente. 10–20 min maximum ; remplace le travail personnel initial.\par\textbf{Note :} Acquisition prévue, à ajuster à partir du travail réellement effectué et des diagnostics. Correction CQFD dans le livret professeur ; reprendre au prochain cours.
\mxneed{8}\section*{S17 -- 2026-11-12 : Quand la dérivabilité échoue}
\textbf{Statut :} cours. 09:45–11:45.\par\textbf{Rituel :} R017, R041.\par\textbf{En classe :} C04-U4 ; exemples résolus et applications ; C04-AP4 puis E7 ou E8.  QCM 2 C04 (10 min) et bilan.\par\textbf{Travail personnel :} C04-E8 à reprendre dans une semaine, 10–20 min maximum.\par\textbf{Note :} Acquisition prévue, à ajuster à partir du travail réellement effectué et des diagnostics.
\mxneed{8}\section*{S18 -- 2026-11-17 : Conditionner : changer le dénominateur}
\textbf{Statut :} cours. 09:45–11:45.\par\textbf{Rituel :} R018, R042.\par\textbf{En classe :} C05-U1 ; exemples résolus et applications ; C05-E1. Diagnostic D : 10 min ; activité RECH : 15 min.\par\textbf{Travail personnel :} C05-E1, 10–20 min maximum.\par\textbf{Note :} Acquisition prévue, à ajuster à partir du travail réellement effectué et des diagnostics.
\mxneed{8}\section*{S19 -- 2026-11-19 : Arbres, partitions et probabilités totales}
\textbf{Statut :} cours. 09:45–11:45.\par\textbf{Rituel :} R019, R043.\par\textbf{En classe :} C05-U2 ; exemples résolus et applications ; C05-E2 et E6. \par\textbf{Travail personnel :} CQFD, exercice 67 p. 336, questions 1, 2 et 3. Construire un arbre en conservant les événements D et N du manuel. Attention : N désigne un test négatif. 10–20 min maximum ; remplace le travail personnel initial.\par\textbf{Note :} Acquisition prévue, à ajuster à partir du travail réellement effectué et des diagnostics. Correction CQFD dans le livret professeur ; reprendre au prochain cours.
\mxneed{8}\section*{S20 -- 2026-11-24 : Indépendance et incompatibilité}
\textbf{Statut :} cours. 09:45–11:45.\par\textbf{Rituel :} R020, R044.\par\textbf{En classe :} C05-U3 ; exemples résolus et applications ; C05-E3.  QCM 1 C05 (10 min, correction comprise).\par\textbf{Travail personnel :} C05-E3, 10–20 min maximum.\par\textbf{Note :} Acquisition prévue, à ajuster à partir du travail réellement effectué et des diagnostics.
\mxneed{8}\section*{S21 -- 2026-11-26 : Épreuves successives et Bernoulli}
\textbf{Statut :} cours. 09:45–11:45.\par\textbf{Rituel :} R021, R045.\par\textbf{En classe :} C05-U4, C05-U5 ; exemples résolus et applications ; C05-E4 et E5 au choix.  QCM 2 C05 (10 min) et bilan.\par\textbf{Travail personnel :} CQFD, exercice 70 p. 336, questions a et b. Dénombrer les quatre issues équiprobables des deux lancers pour réfuter les égalités proposées. 10–20 min maximum ; remplace le travail personnel initial.\par\textbf{Note :} Acquisition prévue, à ajuster à partir du travail réellement effectué et des diagnostics. Correction CQFD dans le livret professeur ; reprendre au prochain cours.
\mxneed{8}\section*{S22 -- 2026-12-01 : Fonction dérivée et fonctions usuelles}
\textbf{Statut :} cours. 09:45–11:45.\par\textbf{Rituel :} R022, R040.\par\textbf{En classe :} C06-U1 ; exemples résolus et applications ; C06-AP1 : dérivées usuelles seulement. Diagnostic D : 10 min ; activité RECH : 15 min.\par\textbf{Travail personnel :} C06-AP1, 10–20 min maximum.\par\textbf{Note :} Acquisition prévue, à ajuster à partir du travail réellement effectué et des diagnostics.
\mxneed{8}\section*{S23 -- 2026-12-03 : Somme, produit, inverse et quotient}
\textbf{Statut :} cours. 09:45–11:45.\par\textbf{Rituel :} R023, R041.\par\textbf{En classe :} C06-U2 ; exemples résolus et applications ; C06-E1. \par\textbf{Travail personnel :} C06-E1, 10–20 min maximum.\par\textbf{Note :} Acquisition prévue, à ajuster à partir du travail réellement effectué et des diagnostics.
\mxneed{8}\section*{S24 -- 2026-12-08 : Signe de la dérivée et extrema}
\textbf{Statut :} cours. 09:45–11:45.\par\textbf{Rituel :} R024, R042.\par\textbf{En classe :} C06-U3 ; exemples résolus et applications ; C06-E2.  QCM 1 C06 (10 min, correction comprise).\par\textbf{Travail personnel :} CQFD, exercice 16 p. 156, questions f et g. Donner pour chaque fonction un intervalle où elle est à la fois négative et croissante ; justifier les deux propriétés séparément. 10–20 min maximum ; remplace le travail personnel initial.\par\textbf{Note :} Acquisition prévue, à ajuster à partir du travail réellement effectué et des diagnostics. Correction CQFD dans le livret professeur ; reprendre au prochain cours.
\mxneed{8}\section*{S25 -- 2026-12-10 : Parité, inégalités et comparaison}
\textbf{Statut :} cours. 09:45–11:45.\par\textbf{Rituel :} R025, R043.\par\textbf{En classe :} C06-U4, C06-U5 ; exemples résolus et applications ; C06-E4 et E5 au choix.  QCM 2 C06 (10 min) et bilan.\par\textbf{Travail personnel :} CQFD, exercice 12 p. 156, questions 1 et 2. Calculer la différence, étudier son signe et traduire celui-ci en position relative des courbes. 10–20 min maximum ; remplace le travail personnel initial.\par\textbf{Note :} Acquisition prévue, à ajuster à partir du travail réellement effectué et des diagnostics. Correction CQFD dans le livret professeur ; reprendre au prochain cours.
\mxneed{8}\section*{S26 -- 2026-12-15 : Reprendre les sommes et les seuils}
\textbf{Statut :} cours. 09:45–11:45.\par\textbf{Rituel :} R026, R038.\par\textbf{En classe :} C01-E4, E6, E7 : contrôler nombre de termes et rang précédent ; reprise produit scalaire C03-E8. Mise à disposition de EV2-T1 : choisir une question adaptée aux notions déjà étudiées.\par\textbf{Travail personnel :} C01-E8, 10 min.\par\textbf{Note :} Acquisition prévue, à ajuster à partir du travail réellement effectué et des diagnostics. Entraînement de réserve : remplace une tâche personnelle prévue, sans cumul de devoirs ; ne pas exiger deux sujets complets à domicile.
\mxneed{8}\section*{S27 -- 2026-12-17 : Consolider dérivation et probabilités}
\textbf{Statut :} cours. 09:45–11:45.\par\textbf{Rituel :} R027, R045.\par\textbf{En classe :} C06-E1 et E2 ; comparer avec C05-E8, deux temps séparés de consolidation. Mise à disposition de EV2-T2 : choisir une question adaptée aux notions déjà étudiées.\par\textbf{Travail personnel :} Revoir les erreurs relevées, 15 min.\par\textbf{Note :} Acquisition prévue, à ajuster à partir du travail réellement effectué et des diagnostics. Entraînement de réserve : remplace une tâche personnelle prévue, sans cumul de devoirs ; ne pas exiger deux sujets complets à domicile.
\mxneed{8}\section*{S28 -- 2027-01-05 : Vecteur normal et équation cartésienne}
\textbf{Statut :} cours. 09:45–11:45.\par\textbf{Rituel :} R028, R052.\par\textbf{En classe :} C07-U1 ; exemples résolus et applications ; C07-E1. Diagnostic D : 10 min ; activité RECH : 15 min.\par\textbf{Travail personnel :} CQFD, exercice 18 p. 279, questions a, b, c et d. Lire un vecteur normal non nul pour chaque équation, sans résoudre en y. 10–20 min maximum ; remplace le travail personnel initial.\par\textbf{Note :} Acquisition prévue, à ajuster à partir du travail réellement effectué et des diagnostics. Correction CQFD dans le livret professeur ; reprendre au prochain cours.
\mxneed{8}\section*{S29 -- 2027-01-07 : Projeté orthogonal et distance}
\textbf{Statut :} cours. 09:45–11:45.\par\textbf{Rituel :} R029, R053.\par\textbf{En classe :} C07-U2 ; exemples résolus et applications ; C07-E2 et E5. \par\textbf{Travail personnel :} CQFD, exercice 21 p. 279, questions a et b. Écrire la condition d’orthogonalité à la normale, puis vérifier que le point donné appartient à la droite. 10–20 min maximum ; remplace le travail personnel initial.\par\textbf{Note :} Acquisition prévue, à ajuster à partir du travail réellement effectué et des diagnostics. Correction CQFD dans le livret professeur ; reprendre au prochain cours.
\mxneed{8}\section*{S30 -- 2027-01-12 : Équation d’un cercle}
\textbf{Statut :} cours. 09:45–11:45.\par\textbf{Rituel :} R030, R054.\par\textbf{En classe :} C07-U3, C07-U4 ; exemples résolus et applications ; C07-E3 et E4.  QCM 1 C07 (10 min, correction comprise). QCM 2 C07 (10 min) et bilan.\par\textbf{Travail personnel :} C07-E3, 10–20 min maximum.\par\textbf{Note :} Acquisition prévue, à ajuster à partir du travail réellement effectué et des diagnostics.
\mxneed{8}\section*{S31 -- 2027-01-14 : Créneau réservé aux examens blancs}
\textbf{Statut :} neutralise. 09:45–11:45.\par\textbf{Rituel :} aucun.\par\textbf{En classe :} Fenêtre examen blanc 1 (14–16 janvier) : placement précis des mathématiques à confirmer.\par\textbf{Travail personnel :} \par\textbf{Note :} Fenêtre examen blanc 1 (14–16 janvier) : placement précis des mathématiques à confirmer.
\mxneed{8}\section*{S32 -- 2027-01-19 : Variable, événement et loi}
\textbf{Statut :} cours. 09:45–11:45.\par\textbf{Rituel :} R001, R056.\par\textbf{En classe :} C08-U1 ; exemples résolus et applications ; C08-E1. Diagnostic D : 10 min ; activité RECH : 15 min.\par\textbf{Travail personnel :} C08-E1, 10–20 min maximum.\par\textbf{Note :} Acquisition prévue, à ajuster à partir du travail réellement effectué et des diagnostics.
\mxneed{8}\section*{S33 -- 2027-01-21 : Espérance et linéarité}
\textbf{Statut :} cours. 09:45–11:45.\par\textbf{Rituel :} R002, R057.\par\textbf{En classe :} C08-U2 ; exemples résolus et applications ; C08-E3 : mise équitable seulement ; variance différée. \par\textbf{Travail personnel :} CQFD, exercice 28 p. 306, questions toutes. Calculer l’espérance à partir du tableau de loi ; conserver les probabilités associées à chaque valeur. 10–20 min maximum ; remplace le travail personnel initial.\par\textbf{Note :} Acquisition prévue, à ajuster à partir du travail réellement effectué et des diagnostics. Correction CQFD dans le livret professeur ; reprendre au prochain cours.
\mxneed{8}\section*{S34 -- 2027-01-26 : Variance, écart type et König–Huygens}
\textbf{Statut :} cours. 09:45–11:45.\par\textbf{Rituel :} R003, R058.\par\textbf{En classe :} C08-U3 ; exemples résolus et applications ; C08-E2.  QCM 1 C08 (10 min, correction comprise).\par\textbf{Travail personnel :} CQFD, exercice 31 p. 306, questions toutes ; utiliser la loi de l’exercice 28. Calculer la variance puis l’écart type, en respectant l’arrondi au dixième demandé par le manuel. 10–20 min maximum ; remplace le travail personnel initial.\par\textbf{Note :} Acquisition prévue, à ajuster à partir du travail réellement effectué et des diagnostics. Correction CQFD dans le livret professeur ; reprendre au prochain cours.
\mxneed{8}\section*{S35 -- 2027-01-28 : Transformer et comparer des gains}
\textbf{Statut :} cours. 09:45–11:45.\par\textbf{Rituel :} R004, R059.\par\textbf{En classe :} C08-U4, C08-U5 ; exemples résolus et applications ; C08-E4 ou E5 puis atelier d’échantillons.  QCM 2 C08 (10 min) et bilan.\par\textbf{Travail personnel :} C08-E8 à reprendre dans une semaine, 10–20 min maximum.\par\textbf{Note :} Acquisition prévue, à ajuster à partir du travail réellement effectué et des diagnostics.
\mxneed{8}\section*{S36 -- 2027-02-02 : Définir une fonction nouvelle}
\textbf{Statut :} cours. 09:45–11:45.\par\textbf{Rituel :} R005, R054.\par\textbf{En classe :} C09-U1 ; exemples résolus et applications ; C09-AP1 : tangente en 1. Diagnostic D : 10 min ; activité RECH : 15 min.\par\textbf{Travail personnel :} C09-AP1, 10–20 min maximum.\par\textbf{Note :} Acquisition prévue, à ajuster à partir du travail réellement effectué et des diagnostics.
\mxneed{8}\section*{S37 -- 2027-02-04 : Propriétés algébriques et comparaisons}
\textbf{Statut :} cours. 09:45–11:45.\par\textbf{Rituel :} R006, R055.\par\textbf{En classe :} C09-U2 ; exemples résolus et applications ; C09-E1 et E2. \par\textbf{Travail personnel :} CQFD, exercice 18 p. 187, questions a, b et c. Simplifier les produits et quotients en citant la propriété des exponentielles employée. 10–20 min maximum ; remplace le travail personnel initial.\par\textbf{Note :} Acquisition prévue, à ajuster à partir du travail réellement effectué et des diagnostics. Correction CQFD dans le livret professeur ; reprendre au prochain cours.
\mxneed{8}\section*{S38 -- 2027-02-09 : Dériver et étudier}
\textbf{Statut :} cours. 09:45–11:45.\par\textbf{Rituel :} R007, R056.\par\textbf{En classe :} C09-U3 ; exemples résolus et applications ; C09-E3.  QCM 1 C09 (10 min, correction comprise).\par\textbf{Travail personnel :} CQFD, exercice 30 p. 187, questions a, b et c. Résoudre les trois inéquations et écrire les ensembles de solutions sous forme d’intervalles. 10–20 min maximum ; remplace le travail personnel initial.\par\textbf{Note :} Acquisition prévue, à ajuster à partir du travail réellement effectué et des diagnostics. Correction CQFD dans le livret professeur ; reprendre au prochain cours.
\mxneed{8}\section*{S39 -- 2027-02-11 : Du discret au continu}
\textbf{Statut :} cours. 09:45–11:45.\par\textbf{Rituel :} R008, R057.\par\textbf{En classe :} C09-U4, C09-U5 ; exemples résolus et applications ; C09-E4 et E5 au choix.  QCM 2 C09 (10 min) et bilan.\par\textbf{Travail personnel :} C09-E8 à reprendre dans une semaine, 10–20 min maximum.\par\textbf{Note :} Acquisition prévue, à ajuster à partir du travail réellement effectué et des diagnostics.
\mxneed{8}\section*{S40 -- 2027-03-02 : Mesurer un angle en radians}
\textbf{Statut :} cours. 09:45–11:45.\par\textbf{Rituel :} R009, R052.\par\textbf{En classe :} C10-U1 ; exemples résolus et applications ; C10-E1. Diagnostic D : 10 min ; activité RECH : 15 min.\par\textbf{Travail personnel :} C10-E1, 10–20 min maximum.\par\textbf{Note :} Acquisition prévue, à ajuster à partir du travail réellement effectué et des diagnostics.
\mxneed{8}\section*{S41 -- 2027-03-04 : Enroulement et image d’un réel}
\textbf{Statut :} cours. 09:45–11:45.\par\textbf{Rituel :} R010, R053.\par\textbf{En classe :} C10-U2 ; exemples résolus et applications ; C10-E2. \par\textbf{Travail personnel :} C10-E2, 10–20 min maximum.\par\textbf{Note :} Acquisition prévue, à ajuster à partir du travail réellement effectué et des diagnostics.
\mxneed{8}\section*{S42 -- 2027-03-09 : Sinus, cosinus et valeurs remarquables}
\textbf{Statut :} cours. 09:45–11:45.\par\textbf{Rituel :} R011, R054.\par\textbf{En classe :} C10-U3 ; exemples résolus et applications ; C10-E3 et E4.  QCM 1 C10 (10 min, correction comprise).\par\textbf{Travail personnel :} CQFD, exercice 39 p. 220, questions 1, 2, 3 et 5. Traiter les quatre questions de calcul et de lecture du cercle ; justifier les réponses sur la copie. 10–20 min maximum ; remplace le travail personnel initial.\par\textbf{Note :} Acquisition prévue, à ajuster à partir du travail réellement effectué et des diagnostics. Correction CQFD dans le livret professeur ; reprendre au prochain cours.
\mxneed{8}\section*{S43 -- 2027-03-11 : Angles associés et symétries}
\textbf{Statut :} cours. 09:45–11:45.\par\textbf{Rituel :} R012, R055.\par\textbf{En classe :} C10-U4, C10-U5 ; exemples résolus et applications ; C10-E6 ou E7.  QCM 2 C10 (10 min) et bilan.\par\textbf{Travail personnel :} CQFD, exercice 32 p. 220, questions a et b. Utiliser la valeur exacte fournie et les angles associés, sans calculatrice. 10–20 min maximum ; remplace le travail personnel initial.\par\textbf{Note :} Acquisition prévue, à ajuster à partir du travail réellement effectué et des diagnostics. Correction CQFD dans le livret professeur ; reprendre au prochain cours.
\mxneed{8}\section*{S44 -- 2027-03-16 : Dérivation et valeur absolue : remédiation}
\textbf{Statut :} cours. 09:45–11:45.\par\textbf{Rituel :} R013, R050.\par\textbf{En classe :} C04-E8 ; C06-E6 ; reprise des preuves du carré et de l’inverse.\par\textbf{Travail personnel :} C06-E8, 10 min.\par\textbf{Note :} Acquisition prévue, à ajuster à partir du travail réellement effectué et des diagnostics.
\mxneed{8}\section*{S45 -- 2027-03-18 : Listes, sommes et simulations}
\textbf{Statut :} cours. 09:45–11:45.\par\textbf{Rituel :} R014, R057.\par\textbf{En classe :} Atelier Python : suites.py et probabilites.py ; C01-E7 et C08-E7, tâches au choix selon besoins.\par\textbf{Travail personnel :} Relire un programme et écrire son résultat à la main, 15 min.\par\textbf{Note :} Acquisition prévue, à ajuster à partir du travail réellement effectué et des diagnostics.
\mxneed{8}\section*{S46 -- 2027-03-23 : Variables aléatoires : consolider les modèles}
\textbf{Statut :} cours. 09:45–11:45.\par\textbf{Rituel :} R015, R064.\par\textbf{En classe :} C08-E2 et E8 ; loi, espérance et variance ; protocole de simulation d’échantillons.\par\textbf{Travail personnel :} C08-E5, 15 min.\par\textbf{Note :} Acquisition prévue, à ajuster à partir du travail réellement effectué et des diagnostics.
\mxneed{8}\section*{S47 -- 2027-03-25 : Préparer le blanc 2 : méthodes}
\textbf{Statut :} cours. 09:45–11:45.\par\textbf{Rituel :} R016, R071.\par\textbf{En classe :} EV3-T1 : exercices, avec une résolution individuelle et une comparaison de méthodes.\par\textbf{Travail personnel :} Finir un exercice désigné du sujet EV3-T1, 20 min.\par\textbf{Note :} Acquisition prévue, à ajuster à partir du travail réellement effectué et des diagnostics.
\mxneed{8}\section*{S48 -- 2027-03-30 : Préparer le blanc 2 : autonomie}
\textbf{Statut :} cours. 09:45–11:45.\par\textbf{Rituel :} R017, R031.\par\textbf{En classe :} EV3-T2 : sélection de deux exercices selon le diagnostic ; verbaliser les contrôles.\par\textbf{Travail personnel :} Revoir la fiche personnelle d’erreurs, 15 min.\par\textbf{Note :} Acquisition prévue, à ajuster à partir du travail réellement effectué et des diagnostics.
\mxneed{8}\section*{S49 -- 2027-04-01 : Créneau réservé aux examens blancs}
\textbf{Statut :} neutralise. 09:45–11:45.\par\textbf{Rituel :} aucun.\par\textbf{En classe :} Fenêtre examen blanc 2 (1–3 avril) : placement précis des mathématiques à confirmer.\par\textbf{Travail personnel :} \par\textbf{Note :} Fenêtre examen blanc 2 (1–3 avril) : placement précis des mathématiques à confirmer.
\mxneed{8}\section*{S50 -- 2027-04-06 : Retour du blanc 2}
\textbf{Statut :} cours. 09:45–11:45.\par\textbf{Rituel :} R019, R045.\par\textbf{En classe :} Corriger les erreurs effectivement observées ; sujets EV3-A/B selon version passée.\par\textbf{Travail personnel :} Refaire une question incorrecte sans consulter le corrigé, 15 min.\par\textbf{Note :} Acquisition prévue, à ajuster à partir du travail réellement effectué et des diagnostics.
\mxneed{8}\section*{S51 -- 2027-04-08 : Consolidation après le blanc 2}
\textbf{Statut :} cours. 09:45–11:45.\par\textbf{Rituel :} R020, R052.\par\textbf{En classe :} C07-E2, E3 et C10-E6 ; deux groupes selon erreurs de géométrie.\par\textbf{Travail personnel :} C07-E8, 10 min.\par\textbf{Note :} Acquisition prévue, à ajuster à partir du travail réellement effectué et des diagnostics.
\mxneed{8}\section*{S52 -- 2027-04-27 : Choisir une méthode en algèbre}
\textbf{Statut :} cours. 09:45–11:45.\par\textbf{Rituel :} R021, R059.\par\textbf{En classe :} C02-E3, E4 et E7\par\textbf{Travail personnel :} Refaire une question ciblée, 10–15 min ; pas de devoir supplémentaire après un sujet complet.\par\textbf{Note :} Acquisition prévue, à ajuster à partir du travail réellement effectué et des diagnostics.
\mxneed{8}\section*{S53 -- 2027-04-29 : Probabilités : identifier le dénominateur}
\textbf{Statut :} cours. 09:45–11:45.\par\textbf{Rituel :} R022, R066.\par\textbf{En classe :} C05-E1, E2 et E6\par\textbf{Travail personnel :} Refaire une question ciblée, 10–15 min ; pas de devoir supplémentaire après un sujet complet.\par\textbf{Note :} Acquisition prévue, à ajuster à partir du travail réellement effectué et des diagnostics.
\mxneed{8}\section*{S54 -- 2027-05-04 : Dérivation : signe et domaine}
\textbf{Statut :} cours. 09:45–11:45.\par\textbf{Rituel :} R023, R073.\par\textbf{En classe :} C06-E1, E3 et E8\par\textbf{Travail personnel :} Refaire une question ciblée, 10–15 min ; pas de devoir supplémentaire après un sujet complet.\par\textbf{Note :} Acquisition prévue, à ajuster à partir du travail réellement effectué et des diagnostics.
\mxneed{8}\section*{S55 -- 2027-05-11 : Géométrie : choisir un cadre}
\textbf{Statut :} cours. 09:45–11:45.\par\textbf{Rituel :} R024, R033.\par\textbf{En classe :} C03-E3 et C07-E2\par\textbf{Travail personnel :} Refaire une question ciblée, 10–15 min ; pas de devoir supplémentaire après un sujet complet.\par\textbf{Note :} Acquisition prévue, à ajuster à partir du travail réellement effectué et des diagnostics.
\mxneed{8}\section*{S56 -- 2027-05-13 : Exponentielle et suites}
\textbf{Statut :} cours. 09:45–11:45.\par\textbf{Rituel :} R025, R040.\par\textbf{En classe :} C09-E2, E3 et E8\par\textbf{Travail personnel :} Refaire une question ciblée, 10–15 min ; pas de devoir supplémentaire après un sujet complet.\par\textbf{Note :} Acquisition prévue, à ajuster à partir du travail réellement effectué et des diagnostics.
\mxneed{8}\section*{S57 -- 2027-05-18 : EAM : entraînement T1 en temps limité}
\textbf{Statut :} cours. 09:45–11:45.\par\textbf{Rituel :} aucun.\par\textbf{En classe :} EV4-T1, sujet complet 120 min\par\textbf{Travail personnel :} Refaire une question ciblée, 10–15 min ; pas de devoir supplémentaire après un sujet complet.\par\textbf{Note :} Acquisition prévue, à ajuster à partir du travail réellement effectué et des diagnostics.
\mxneed{8}\section*{S58 -- 2027-05-20 : Retour de l’entraînement T1}
\textbf{Statut :} cours. 09:45–11:45.\par\textbf{Rituel :} R027, R054.\par\textbf{En classe :} EV4-T1 : correction différenciée ; refaire deux questions\par\textbf{Travail personnel :} Refaire une question ciblée, 10–15 min ; pas de devoir supplémentaire après un sujet complet.\par\textbf{Note :} Acquisition prévue, à ajuster à partir du travail réellement effectué et des diagnostics.
\mxneed{8}\section*{S59 -- 2027-05-25 : Variance et statistiques}
\textbf{Statut :} cours. 09:45–11:45.\par\textbf{Rituel :} R028, R061.\par\textbf{En classe :} C08-E2 et E4 ; R017–R020 et R026–R028\par\textbf{Travail personnel :} Refaire une question ciblée, 10–15 min ; pas de devoir supplémentaire après un sujet complet.\par\textbf{Note :} Acquisition prévue, à ajuster à partir du travail réellement effectué et des diagnostics.
\mxneed{8}\section*{S60 -- 2027-05-27 : Cercle, radians et preuves}
\textbf{Statut :} cours. 09:45–11:45.\par\textbf{Rituel :} R029, R068.\par\textbf{En classe :} C10-E2, E3 et E6 ; démonstrations des valeurs remarquables\par\textbf{Travail personnel :} Refaire une question ciblée, 10–15 min ; pas de devoir supplémentaire après un sujet complet.\par\textbf{Note :} Acquisition prévue, à ajuster à partir du travail réellement effectué et des diagnostics.
\mxneed{8}\section*{S61 -- 2027-06-01 : EAM : entraînement T2 en temps limité}
\textbf{Statut :} cours. 09:45–11:45.\par\textbf{Rituel :} aucun.\par\textbf{En classe :} EV4-T2, sujet complet 120 min\par\textbf{Travail personnel :} Refaire une question ciblée, 10–15 min ; pas de devoir supplémentaire après un sujet complet.\par\textbf{Note :} Acquisition prévue, à ajuster à partir du travail réellement effectué et des diagnostics.
\mxneed{8}\section*{S62 -- 2027-06-03 : Retour de l’entraînement T2}
\textbf{Statut :} cours. 09:45–11:45.\par\textbf{Rituel :} R078, R035.\par\textbf{En classe :} EV4-T2 : correction différenciée ; plan individuel\par\textbf{Travail personnel :} Refaire une question ciblée, 10–15 min ; pas de devoir supplémentaire après un sujet complet.\par\textbf{Note :} Acquisition prévue, à ajuster à partir du travail réellement effectué et des diagnostics.
\mxneed{8}\section*{S63 -- 2027-06-08 : Consolider les preuves}
\textbf{Statut :} cours. 09:45–11:45.\par\textbf{Rituel :} R001, R042.\par\textbf{En classe :} TRANS : logique ; C01-U4, C03-U4 et C06-U2\par\textbf{Travail personnel :} Refaire une question ciblée, 10–15 min ; pas de devoir supplémentaire après un sujet complet.\par\textbf{Note :} Acquisition prévue, à ajuster à partir du travail réellement effectué et des diagnostics.
\mxneed{8}\section*{S64 -- 2027-06-10 : Synthèse algèbre-analyse}
\textbf{Statut :} cours. 09:45–11:45.\par\textbf{Rituel :} R002, R049.\par\textbf{En classe :} C06-E5 et C09-E6\par\textbf{Travail personnel :} Refaire une question ciblée, 10–15 min ; pas de devoir supplémentaire après un sujet complet.\par\textbf{Note :} Acquisition prévue, à ajuster à partir du travail réellement effectué et des diagnostics.
\mxneed{8}\section*{S65 -- 2027-06-15 : Automatismes et lecture de données}
\textbf{Statut :} cours. 09:45–11:45.\par\textbf{Rituel :} R003, R056.\par\textbf{En classe :} R001–R030, sélection selon erreurs ; R078\par\textbf{Travail personnel :} Refaire une question ciblée, 10–15 min ; pas de devoir supplémentaire après un sujet complet.\par\textbf{Note :} Acquisition prévue, à ajuster à partir du travail réellement effectué et des diagnostics.
\mxneed{8}\section*{S66 -- 2027-06-17 : Dernières reprises et méthode d’épreuve}
\textbf{Statut :} cours. 09:45–11:45.\par\textbf{Rituel :} R004, R063.\par\textbf{En classe :} Questions déjà ratées ; gestion des 120 minutes, vérification des domaines\par\textbf{Travail personnel :} Refaire une question ciblée, 10–15 min ; pas de devoir supplémentaire après un sujet complet.\par\textbf{Note :} Acquisition prévue, à ajuster à partir du travail réellement effectué et des diagnostics.
\mxneed{8}\section*{S67 -- 2027-06-22 : Consolidation selon maintien des cours}
\textbf{Statut :} conditionnel. 09:45–11:45.\par\textbf{Rituel :} R005, R070.\par\textbf{En classe :} Reprises ciblées ou transition vers la terminale : C06-E8, C09-E8, C10-E8, selon présence et examens.\par\textbf{Travail personnel :} \par\textbf{Note :} Ne pas compter ce créneau dans la couverture obligatoire avant l’EAM. Maintien des cours à confirmer.
\mxneed{8}\section*{S68 -- 2027-06-24 : Consolidation selon maintien des cours}
\textbf{Statut :} conditionnel. 09:45–11:45.\par\textbf{Rituel :} R006, R077.\par\textbf{En classe :} Reprises ciblées ou transition vers la terminale : C06-E8, C09-E8, C10-E8, selon présence et examens.\par\textbf{Travail personnel :} \par\textbf{Note :} Ne pas compter ce créneau dans la couverture obligatoire avant l’EAM. Maintien des cours à confirmer.
\mxneed{8}\section*{S69 -- 2027-06-29 : Consolidation selon maintien des cours}
\textbf{Statut :} conditionnel. 09:45–11:45.\par\textbf{Rituel :} R007, R037.\par\textbf{En classe :} Reprises ciblées ou transition vers la terminale : C06-E8, C09-E8, C10-E8, selon présence et examens.\par\textbf{Travail personnel :} \par\textbf{Note :} Ne pas compter ce créneau dans la couverture obligatoire avant l’EAM. Maintien des cours à confirmer.
\mxneed{8}\section*{S70 -- 2027-07-01 : Consolidation selon maintien des cours}
\textbf{Statut :} conditionnel. 09:45–11:45.\par\textbf{Rituel :} R008, R044.\par\textbf{En classe :} Reprises ciblées ou transition vers la terminale : C06-E8, C09-E8, C10-E8, selon présence et examens.\par\textbf{Travail personnel :} \par\textbf{Note :} Ne pas compter ce créneau dans la couverture obligatoire avant l’EAM. Maintien des cours à confirmer.
